% Auxiliary exact facts; the general product step remains unproved.
\section{Which arm properties require increasing parts?}
For an arbitrary finite list of positive integers $T$, define the
formal arm polynomial by $A_\varnothing=1$ and
$A_{Tt}=1+S_tA_T$. This agrees with the principal-ideal polynomial
when $T$ is increasing. Without that hypothesis it need not count
a composition ideal, and it will not be identified with one below.

The ordinary unimodality assertion for this formal recurrence is
already present, after reversing the list convention, in
B.~E.~Sagan, \emph{Compositions inside a rectangle and unimodality},
J.~Algebraic Combin.~29 (2009), 405--411, Theorem~3.3;
the author's May 5, 2008 manuscript was checked at
\url{https://users.math.msu.edu/users/sagan/papers/old/cir.pdf}.
That theorem concerns his coordinatewise composition order, not
generalized subword order. It is background here, not a resolution
of the unrestricted conjecture. The additional rising-difference
and kernel comparisons below are stated separately.

\begin{lemma}[Uniform windows and rising first differences]
Let $p$ be nonnegative and unimodal, and suppose the positive
part of its first-difference sequence is unimodal. Then the same
is true of $[t]p$ for each $t\ge1$.
\end{lemma}
\begin{proof}
Let $m$ be the last mode of $p$ and extend all sequences by zero.
For $n\le m$, the difference of $q=[t]p$ is a length-$t$ window
of the nonnegative rising differences of $p$. These differences,
extended by zero, are unimodal, so their window is unimodal.
For $m\le n\le m+t-1$, the differences of $q$ satisfy
\[
 \Delta q_{n+1}-\Delta q_n
  =(p_{n+1}-p_n)-(p_{n-t+1}-p_{n-t})\le0.
\]
For $n\ge m+t$, $\Delta q_n=p_n-p_{n-t}\le0$.
Thus the positive first differences form a prefix of a unimodal
window followed by a nonincreasing portion and then zeros.
This proves the claim, including the joining comparison at $m$.
\end{proof}

\begin{proposition}
For every formal arm $T$, $A_T$ is unimodal. Its reverse is
log-concave on its rising side, and the positive part of the
first differences of that reverse is unimodal. The same two
rising-side properties hold for $\delta_m(Y)$ when $Y$ is
increasing and $m\ge\max(1,\kappa(Y)-1)$.
\end{proposition}
\begin{proof}
For the formal arm, its reverse satisfies
$q_{Tt}=[t]q_T+z^{\operatorname{wt}T+t}$, for any order of the
parts. The convolution has final coefficient one; appending one
does not change the positive part of its first differences.
Use the uniform-window lemma above and the earlier rising-side
log-concavity window lemma inductively.

For a strict kernel, the case $b<M$ has rising differences equal
to a prefix of $q_Y$, which is unimodal. In the case $b\ge M$,
the added interval in $[t]\delta_M(Z)+z^b[M-t]$ affects only
the descending portion and does not introduce positive differences.
The window lemma applies. At $t=m$ the kernel recurrence is a
window without correction. At $t>m$, the subcritical-kernel proof
shows that subtraction affects only the descending portion and
that the ending splice is nonincreasing. The same induction
therefore applies. Rising-side log-concavity was proved separately
in the subcritical-kernel theorem.
\end{proof}

\begin{proposition}[Pure-arm ratio comparisons do not require sorted parts]
Let every part of $T$ be at least two. Put
$R_T=\prod_{t\in T}S_t$, $K_T=(1-z)A_T+zR_T$,
$u=\operatorname{wt}T$, and $T_T=\prod_{t\in T}[t]$.
Then $K_T$ is nonnegative, and for each $\alpha\ge0$ the sequences
\[
 a_i=[z^i]z^uK_T(1/z),\qquad b_i=[z^{i+\alpha+1}]T_T
\]
satisfy $a_jb_i-a_ib_j\ge0$ whenever $i<j$.
\end{proposition}
\begin{proof}
The identity $K_{Tt}=1-z+S_tK_T$ holds formally. Starting from
$K_\varnothing=1$, the constant coefficient stays one and the
linear coefficient stays zero, so all coefficients remain
nonnegative. In reverse the recurrence is
\[
 K_{Tt}^*=[t]K_T^*-z^{u+t-1}+z^{u+t},\qquad
 T_{Tt}=[t]T_T.
\]
The proof of the pure-arm ratio theorem uses only these identities,
nonnegativity of $K$, and the degree $u-\operatorname{len}T$ of
$T_T$. The negative endpoint correction lies strictly to the
right of the shifted denominator support. Uniform convolution
preserves the ratio comparisons. Larger shifts follow from
log-concavity of $T_T$; its positive support is an interval.
None of these steps uses an ordering of the parts.
\end{proof}

\paragraph{An exact initial exchange.}
For $\alpha\ge1$ and $t\ge2$,
\[
 A_{1^\alpha t}=1+z[t][\alpha+1]
              =A_{1^{t-1}(\alpha+1)}.
\]
Any identical formal suffix preserves this identity. In particular,
when $\alpha+1\le t$ and the remaining parts are at least $t$,
the exchange gives an equality of genuine increasing-arm rank
polynomials. When $\alpha=0$, the first part $t$ can instead be
replaced by $t$ leading ones.

\begin{proposition}[Forward rising differences after the initial ones]
Let $T=1^\alpha U$, where $\alpha\ge0$ and every part of the
formal list $U$ is at least two; the parts of $U$ need not be sorted.
For $p=A_T$, the sequence
\[
 \bigl(\max(p_n-p_{n-1},0)\bigr)_{n\ge1}
\]
is unimodal. If $p$ is nonrectangular, it increases strictly at
every rank from two through its first mode.
\end{proposition}
\begin{proof}
An interval polynomial is immediate. Its first nonrectangular
extension is $1+z[t][\alpha+1]$, whose positive differences after
rank zero form an interval of ones beginning at rank two.
Proceed by induction under appending $t\ge2$. Write $m$ for
the first mode of the old nonrectangular polynomial $p$.
Its positive differences have support exactly $[2,m]$: this
follows from their unimodality, the value $p_2-p_1=1$, and
the definition of the first mode. They are positive integers
throughout this support.

For $q=1+S_tp$ one has $\Delta q_1=0$ and
\[
 \Delta q_n=p_{n-1}-p_{n-t-1}\qquad(n\ge2).
\]
If $t>m$, the segment $2\le n\le t$ is a prefix of the
unimodal sequence $p_{n-1}$. At $n=t+1$ the difference cannot
increase, because $p_t\le p_{t-1}$ and the new subtracted
coefficient is $p_0=1$. From there until $m+t+1$, the upper
index lies on the descending side and the lower index on the
rising side, so the differences are nonincreasing. Beyond that
range they are nonpositive. This gives the assertion in this case.

If $t\le m$, the initial segment $2\le n\le t$ is
nondecreasing. At its junction with the next segment,
\[
 \Delta q_{t+1}-\Delta q_t=(p_t-p_{t-1})-1\ge0.
\]
For $t+1\le n\le m+1$, the differences are the length-$t$
window of the old positive differences (zero-extended outside
$[2,m]$). This window is unimodal. Its values through $n=t+1$
are cumulative sums, so joining it to the preceding increasing
segment cannot create a valley. From $n=m+1$ through $m+t+1$,
the crossing-of-the-mode argument again makes the differences
nonincreasing, and afterward they are nonpositive. The positive
part is therefore unimodal in this case as well.

Finally $q_2-q_1=1$. Its unimodal positive-difference sequence
cannot have an internal zero before reaching the first mode,
which proves the strict-increase assertion and closes induction.
\end{proof}

\paragraph{Scope of these extensions.}
The properties of the reverse kernel still use the increasing
condition on $Y$. A formal arm with unsorted parts does not in
general have a unimodal relative kernel. For example, repeated
terminal parts equal to one can create gaps in the positive
support of its growth deficit. Neither the pure-arm extension
nor the first-difference property is a proof of convolution
unimodality: the same sequence $(1,3,2,1,1,1,1)$ used earlier
also has unimodal positive first differences, but its square
is not unimodal.

\begin{proposition}[Ordinary products with no later return to one]
If $T=1^\alpha U$ and $W=1^\beta V$, with all parts of the
formal skeletons $U,V$ at least two, then $A_TA_W$ is unimodal.
The skeletons may have arbitrary internal order. Its reverse
also has unimodal positive first differences and is log-concave
on its rising side.
\end{proposition}
\begin{proof}
For every formal arm, the lowest height layer is $[0,a]$,
and every higher layer has endpoint sum at least $a+1$.
This follows by induction from $1+S_tA$: uniform convolution
gives all layers of $S_tA$ endpoint sum at least $a+t+1$,
and adding the initial constant only changes its lowest layer.
In the stated class, every nonrectangular polynomial has second
layer exactly $[2,a-1]$, since both coefficients at those two
ranks are at least two and the adjacent end coefficients are one.

The ordinary-product proof therefore applies verbatim within
this deletion-closed class. For completeness, write $a\ge b$.
Layer centers give growth through $\lceil(a+b)/2\rceil$.
When $a=b$ and both factors are nonrectangular, the product of
the two second layers pays for the only negative unit from the
two lowest layers at rank $a+1$, extending growth to that rank.
Deleting a last part $t$ of the heavier arm gives
$A_TA_W=A_W+S_t(A_{T^-}A_W)$. By induction the window term
is unimodal. If $a>b$, its tail joins the guaranteed increasing
segment after the correction $A_W$ has vanished; if $a=b$,
use the extra rank of growth just proved. Rectangular factors
are uniform-window base cases. Thus the product is unimodal.

The reverse-prefix proof in \texttt{joint\_rising\_differences.tex}
uses only this product unimodality, the same strengthened midpoint,
and the formal append recurrence. Those inputs all hold within
the present deletion-closed class, proving the two additional
reverse properties.
\end{proof}

The corresponding formal kernel need not be unimodal even when
there are no ones at all. For $M=5$ and the unsorted formal list
$T=(4,4,4,4,2,2)$, its kernel has coefficients $2,3,2,3$ at
ranks $4,5,6,7$. Thus the increasing hypothesis on the kernel
arm remains essential to the arguments being used here.
